\documentclass[10pt,a4paper]{amsart}
\usepackage[margin=19mm]{geometry}
\usepackage{amsmath,amssymb,mathtools}
\usepackage[scr=rsfs,cal=euler]{mathalfa}
\usepackage{xfrac}
\usepackage[unicode,hidelinks,bookmarks=false]{hyperref}
\newcommand{\ie}{i.e.\ }
\newcommand{\cf}{cf.\ }
\newcommand{\resp}{respectively\ }
\newcommand{\Subsection}{\subsection}
\providecommand{\nobreakdash}{\nobreak\hbox{-}}
\setlength{\emergencystretch}{2em}
\multlinegap0pt
\usepackage{amssymb}
\usepackage{xcolor}
\usepackage{mathtools}
\usepackage[shortlabels]{enumitem}
\usepackage{float}
\usepackage{tikz}
\newtheorem{lemma}{Lemma}
\usepackage{cleveref}
\crefname{lemma}{Lemma}{Lemmas}
\title{Silver Rate Is (Almost) Optimal for Gradient Descent}
\author{Yuhan Ye, Kaizhao Liu}
\date{Source: arXiv:2609.09152v2 (CC BY 4.0)}
\begin{document}
\maketitle

\noindent\emph{Source and licence.} Silver Rate Is (Almost) Optimal for Gradient Descent. Version arXiv:2609.09152v2 (CC BY 4.0), distributed by arXiv under the Creative Commons Attribution 4.0 International licence, http://creativecommons.org/licenses/by/4.0/, as recorded in the arXiv licence metadata for this exact version and shown on the abstract page https://arxiv.org/abs/2609.09152v2. Full licence notice: \texttt{licenses/arxiv-2609.09152-CC-BY-4.0.txt}.\par\smallskip

\noindent\textit{Editorial scope.} Counted unit: the article's lemma lem:weighted-splitting (source0.tex lines 1564--1575). The article's complete original proof of it is delivered with it (source0.tex lines 1587--1715, one \texttt{proof} environment of 129 lines, which the article places away from the statement). No mathematics is written, repaired or re-derived: the statement and the proof are the article's own lines, in the article's own notation, and no retained line is substituted or re-flowed.  Retained uncounted as the source-local context the proof needs: lines 573--576; lines 1350--1360; lines 1515--1522; lines 1581--1585.  Nothing was omitted from any retained span, so the package carries no omission record.  No retained line contains a citation, so the wrapper carries no bibliography and no bibliography entry.  No retained line contains a figure environment, an image-inclusion command or drawing code, so no image is delivered and the package declares no asset dependency.  Editorial formatting only, in addition: an AMS \texttt{amsart} preamble that restates the article's own declarations and macros used by the retained lines, namely the environments \texttt{lemma} on separate counters, together with the standard packages those lines need.  The retained environments print in this document as Lemma 1 in order of appearance, whereas the article prints the same items with the numbering of its own full document.  The complete native source of the article as distributed in the arXiv e-print for this exact version is bundled unchanged as \texttt{sources/209783/source0.tex}.
\begin{equation}
  w^\nu+\kappa B^\nu\le x^\nu+y^\nu.
  \label{global:eq-scalar-hypothesis}
\end{equation}

\begin{lemma}[Log-submodularity]
\label{global:lem-log-submodularity}
For every positive schedule $h\in(0,\infty)^n$ and every two checkpoint sets $A,B$,
\begin{equation}
  \Psi_\lambda(A;h)\Psi_\lambda(B;h)
  \ge\Psi_\lambda(A\cap B;h)\Psi_\lambda(A\cup B;h).
  \label{global:eq-log-submodularity}
\end{equation}
The inequality is strict when $A$ and $B$ are disjoint and nonempty.
Consequently, tight sets are closed under unions and intersections, and two nonempty tight sets cannot be disjoint.
\end{lemma}

\begin{align}
  \widetilde\Psi_\Lambda(S;g,B)
    &:=Y(S)\prod_{j=1}^m\frac{(c-\epsilon)(a+L_j(S))}{B_{i_j}},
    \label{global:eq-pure-cost}\\
  \widetilde\Psi_\Lambda(\varnothing;g,B)
    &:=\Lambda+\sum_{i=1}^k(g_i+B_i).
    \label{global:eq-pure-empty}
\end{align}

\begin{equation}
  J(B):=W^\nu+\kappa\sum_{i=1}^kB_i^\nu,
  \qquad W=\Lambda+\sum_{i=1}^k(g_i+B_i),
  \label{global:eq-weighted-objective}
\end{equation}

\begin{lemma}[Splitting with a correction term]
\label{lem:weighted-splitting}
Fix $\nu\in[1/2,1)$ and $\kappa\ge0$.
Suppose that, for every $x,y\ge a$ and $B>0$, the relations $w=x+y+B-a$ and $wB=(c-\epsilon) xy$ imply \eqref{global:eq-scalar-hypothesis}.
For any auxiliary problem \eqref{global:eq-pure-cost}--\eqref{global:eq-pure-empty} in which the empty set is a minimizer, write $W=\widetilde\Psi_\Lambda(\varnothing;g,B)$.
Then
\begin{equation}
  W^\nu+\kappa\sum_{i=1}^kB_i^\nu
  \le\Lambda^\nu+\sum_{i=1}^k(a+g_i)^\nu.
  \label{global:eq-weighted-splitting}
\end{equation}
\end{lemma}

\begin{proof}[Proof of \cref{lem:weighted-splitting}]
We induct on $k$.
For $k=0$, $W=\Lambda$ and the assertion is equality.
For $k\ge1$, hold $g$ and $\Lambda$ fixed and maximize $J(B)$ from \eqref{global:eq-weighted-objective} over $B_i\ge0$ for which the empty set minimizes $\widetilde\Psi_\Lambda$.
It suffices to prove the desired bound at a maximizer, since its right side is fixed.

\paragraph{Compactness and zero coordinates.}
The feasible set is nonempty because $B=0$ is feasible.
For the singleton constraint at $i$, define
\[
  x_i=a+\sum_{r\le i}g_r+\sum_{r<i}B_r,
  \qquad y_i=\Lambda+\sum_{r>i}(g_r+B_r).
\]
Its constraint is $WB_i\le(c-\epsilon) x_i y_i$.
Since $W\ge y_i>0$, we have $B_i\le(c-\epsilon) x_i$.
The latter bound depends only on earlier variables and fixed values of $g_i$, so applying it successively bounds every $B_i$.
For each nonempty $S$, multiplying $W\le\widetilde\Psi_\Lambda(S;g,B)$ by $\prod_{i\in S}B_i$, when all selected sizes are positive, gives a continuous polynomial inequality.
If a selected size is zero, the left side becomes zero and the right side is positive, agreeing with the $+\infty$ convention for $\widetilde\Psi_\Lambda$.
The feasible set is therefore closed and bounded, and the continuous function $J$ attains a maximum.

If a maximizing vector has $B_i=0$ with $i<k$, delete that checkpoint and replace $g_{i+1}$ by $g_i+g_{i+1}$.
This preserves the values of $\widetilde\Psi_\Lambda$ for all sets that omit $i$, and hence preserves empty-set optimality.
The induction hypothesis, followed by
\[
  (a+g_i+g_{i+1})^\nu
  \le(a+g_i)^\nu+(a+g_{i+1})^\nu,
\]
gives the desired bound.
If $B_k=0$, delete it and replace $\Lambda$ by $\Lambda+g_k$. Then use $(\Lambda+g_k)^\nu\le\Lambda^\nu+(a+g_k)^\nu$.
We may therefore assume that a maximizing vector has every $B_i>0$.

\paragraph{Finding a tight set with one checkpoint.}
At least one nonempty subset is tight.
Otherwise, a sufficiently small increase of any $B_i$ would preserve all constraints and strictly increase $J$.
The auxiliary expression $\widetilde\Psi_\Lambda$ has the same log-submodularity and strictness property as in \cref{global:lem-log-submodularity}. Its insertion ratio is
\[
  \frac{(c-\epsilon)(a+\ell)y}{B(\ell+B+y)},
\]
which is strictly increasing in $\ell$ and $y$ for $y\ge a$ and $B>0$.
The proof by successive insertions applies unchanged.
Choose an inclusion-minimal nonempty tight set $T_0$.
Every nonempty tight set intersects $T_0$, and that intersection is tight.
Minimality therefore implies that $T_0$ lies in every nonempty tight set.

Suppose first that $T_0$ contains two unequal sizes $u<v$, and put $r=v/u>1$.
Because both checkpoints are selected in every nonempty tight set, $\widetilde\Psi_\Lambda$ for each such set depends on them only through the reciprocal product $1/(uv)$.
Define
\[
  A_*:=\frac{v(W+u)}{u(W+v)},\qquad
  D_*:=\frac{W^{\nu-1}+\kappa u^{\nu-1}}
             {W^{\nu-1}+\kappa v^{\nu-1}}.
\]
The positive term $\Lambda$ gives $W>u+v$.
Since $\nu\ge1/2$ and $\kappa\ge0$,
\begin{equation*}
  A_*>\frac{r(r+2)}{2r+1}\ge\sqrt r
  \ge r^{1-\nu}\ge D_*\ge1.
\end{equation*}
For completeness, the first inequality follows because $r(z+1)/(z+r)$ increases in $z=W/u>1+r$.
The second follows, on writing $q=\sqrt r$, from $q^3-2q^2+2q-1=(q-1)(q^2-q+1)\ge0$.
The bound on $D_*$ follows by adding the same positive number $W^{\nu-1}$ to the numerator and denominator of a ratio at most $u^{\nu-1}/v^{\nu-1}=r^{1-\nu}$.

Choose $\theta\in(D_*,A_*)$ and perturb $u\mapsto u-t$, $v\mapsto v+\theta t$.
For every nonempty tight set, $\widetilde\Psi_\Lambda$ initially equals $W$.
The derivative of the difference between its value and the empty-set value is
\[
  W\left(\frac1u-\frac\theta v\right)-(\theta-1)>0,
\]
where the inequality is equivalent to $\theta<A_*$.
Meanwhile the objective has derivative
\[
  \nu\bigl[(\theta-1)W^{\nu-1}
     +\kappa(\theta v^{\nu-1}-u^{\nu-1})\bigr]>0,
\]
since $\theta>D_*$.
Positivity of the $B_i$ and the strictness of all other constraints persist for sufficiently small $t>0$.
This contradicts maximality.

If $T_0$ has at least two elements and all of its sizes are equal to $b$, perturb two of them to $b-t$ and $b+t+t^2/(2b)$.
Their product becomes $b^2-t^2/2-t^3/(2b)$, while $W$ increases by $t^2/(2b)$.
Thus every nonempty tight constraint gains slack
\[
  \frac{W-b}{2b^2}t^2+O(t^3)>0.
\]
The change in $J$ is
\[
  \frac{\nu t^2}{2b}
  \bigl[W^{\nu-1}+\kappa(2\nu-1)b^{\nu-1}\bigr]+O(t^3)>0.
\]
The leading coefficient is positive even for $\nu=1/2$ or $\kappa=0$, because $W^{\nu-1}>0$.
Again the remaining constraints retain their strict slack for small $t$.
This is a contradiction.
It follows that $T_0=\{r\}$ for some index $r$.

\paragraph{Splitting at that checkpoint.}
Write
\[
  x:=a+\sum_{i\le r}g_i+\sum_{i<r}B_i,
  \qquad y:=\Lambda+\sum_{i>r}(g_i+B_i).
\]
Tightness of $\{r\}$ gives $W=x+y+B_r-a$ and $WB_r=(c-\epsilon) xy$, with $x,y\ge a$.
Define the left and right auxiliary expressions by
\[
  \widetilde\Psi_L(S_L):=\widetilde\Psi_{a+g_r}(S_L;g_{1:r-1},B_{1:r-1}),
  \qquad
  \widetilde\Psi_R(S_R):=\widetilde\Psi_\Lambda(S_R;g_{r+1:k},B_{r+1:k}),
\]
where $S_L\subseteq\{1,\ldots,r-1\}$ and $S_R\subseteq\{r+1,\ldots,k\}$ use the original indices; in $\widetilde\Psi_R(S_R)$, each index is reduced by $r$.
For every such pair,
\begin{equation}
  \widetilde\Psi_\Lambda(S_L\cup\{r\}\cup S_R;g,B)
  =\frac{c-\epsilon}{B_r}\widetilde\Psi_L(S_L)\widetilde\Psi_R(S_R).
  \label{global:eq-pure-factorization}
\end{equation}
Their empty-set values are $\widetilde\Psi_L(\varnothing)=x$ and $\widetilde\Psi_R(\varnothing)=y$.
Taking $S_R=\varnothing$ in \eqref{global:eq-pure-factorization} and using optimality of the empty set in the full auxiliary problem gives $((c-\epsilon) y/B_r)\widetilde\Psi_L(S_L)\ge W=((c-\epsilon) y/B_r)x$.
Thus the empty set minimizes the left problem.
Taking $S_L=\varnothing$ gives the same conclusion for the right problem.

The induction hypothesis applies to both sides, including a side with no checkpoints, and yields
\begin{align*}
  x^\nu+\kappa\sum_{i<r}B_i^\nu
    &\le\sum_{i\le r}(a+g_i)^\nu,\\
  y^\nu+\kappa\sum_{i>r}B_i^\nu
    &\le\Lambda^\nu+\sum_{i>r}(a+g_i)^\nu.
\end{align*}
The scalar hypothesis gives $W^\nu+\kappa B_r^\nu\le x^\nu+y^\nu$.
Combining these three inequalities proves \eqref{global:eq-weighted-splitting}.
\end{proof}

\end{document}
